% !PRINTED 199
\weberpage{199}
%<<CONT
\gesperrt{nur nicht gerade der identischen, eine Multiplication
entspricht. Die transformirende Substitution $L$ ist dadurch
selbst bis auf eine Multiplication bestimmt.}

\begin{center}
\S.\,52.

\gesperrt{Die verschiedenen Arten möglicher Gruppen.}
\end{center}

Es sei $a$ einer der Pole der Gruppe $G$, und wir wollen
annehmen, es gebe $\nu - 1$ und nicht mehr Elemente in $G$,
$\Theta_1, \Theta_2, \ldots, \Theta_{\nu-1}$, so dass
\weq{1}{a = \Theta_1(a) = \Theta_2(a) \cdots = \Theta_{\nu-1}(a)}
ist. Ein solcher Pol soll ein $\nu$-zähliger Pol heissen. Es ist
nun zunächst klar, dass die Elemente
\weq{2}{1, \Theta_1, \Theta_2, \ldots, \Theta_{\nu-1}}
für sich eine Gruppe, und zwar einen Theiler von $G$ bilden;
denn aus
\[
a = \Theta_1(a), \quad a = \Theta_2(a)
\]
folgt, wenn man auf der rechten Seite der zweiten Gleichung
$a$ durch das ihm gleiche $\Theta_1(a)$ ersetzt:
\[
a = \Theta_2 \Theta_1(a).
\]
\gesperrt{Es muss also $\nu$ ein Theiler von $n$ sein:}
\weq{3}{n = \nu\mu.}
Wir bezeichnen die Gruppe (2) mit $Q$.

Es lässt sich beweisen, dass diese Gruppe cyklisch sein
muss, also aus den Wiederholungen eines ihrer Elemente
besteht.

Denn wenn wir durch Transformation nach dem Satze 2.
$a$ nach unendlich werfen, so erhalten die Substitutionen (2) die
Form:
\[
x, \; \varepsilon_1 x + c_1, \; \varepsilon_2 x + c_2, \ldots, \varepsilon_{\nu-1} x + c_{\nu-1},
\]
und die Composition von zweien unter ihnen giebt:
\[
\begin{aligned}
\Theta_1 \Theta_2 &= \varepsilon_1 \varepsilon_2 x + (c_2 \varepsilon_1 + c_1) \\
\Theta^\lambda &= \varepsilon^\lambda x + c(\varepsilon^{\lambda-1} + \varepsilon^{\lambda-2} + \cdots + 1) \\
\Theta^{-1} &= \varepsilon^{-1} x - c \varepsilon^{-1}.
\end{aligned}
\]
Daraus schliesst man, dass alle $\varepsilon_1, \varepsilon_2, \ldots, \varepsilon_{\nu-1}$ Einheitswurzeln
vom Grade $\nu$ sein müssen. Ausserdem können nicht zwei
der $\varepsilon$ einander gleich sein. Denn wäre z. B. $\varepsilon_1 = \varepsilon_2$, so wäre
\[
\Theta_1 \Theta_2^{-1} = x + (c_1 - c_2).
\]
